% ===== PRÉAMBULE CANONIQUE — à coller VERBATIM en tête de CHAQUE .tex =====
\documentclass[11pt,a4paper]{article}
\usepackage[utf8]{inputenc}\usepackage[T1]{fontenc}\usepackage{lmodern}
\usepackage[french]{babel}
\usepackage{amsmath,amssymb,mathtools}\usepackage{icomma}
\usepackage{siunitx}\sisetup{locale=FR}  % \qty{}{}, \si{}, \ang{}
\usepackage{xcolor}\usepackage{colortbl}
\usepackage{tikz}\usepackage{pgfplots}\pgfplotsset{compat=1.18}
\usepackage[siunitx,europeanresistors]{circuitikz} % circuits (élec.)
\usepackage[most]{tcolorbox}\usepackage{enumitem}\usepackage{array,booktabs}
\usepackage[a4paper,margin=2.2cm]{geometry}
\usetikzlibrary{arrows.meta,calc,angles,quotes,positioning,patterns,%
  backgrounds,shapes.geometric,decorations.pathmorphing,decorations.markings,intersections}
\definecolor{primary}{RGB}{222,49,99}\definecolor{primarydark}{RGB}{150,22,55}
\definecolor{defcol}{RGB}{0,114,178}\definecolor{methcol}{RGB}{52,92,148}
\definecolor{excol}{RGB}{0,135,90}\definecolor{attncol}{RGB}{200,40,55}
\tcbset{boxbase/.style={enhanced,breakable,boxrule=0.9pt,arc=2.5pt,
  left=3mm,right=3mm,top=2.3mm,bottom=2.3mm,fonttitle=\bfseries,coltitle=white}}
% --- boîtes TUTORIEL (noms EXACTS, ne pas renommer) ---
\newtcolorbox{cle}[1][]{boxbase,colback=defcol!6,colframe=defcol,
  title={\textsc{À retenir}\ifx\relax#1\relax\relax\else\textmd{ : \itshape #1}\fi}}
\newtcolorbox{methode}[1][]{boxbase,colback=methcol!7,colframe=methcol,sharp corners,
  title={\textsc{Méthode}\ifx\relax#1\relax\relax\else\textmd{ --- \itshape #1}\fi}}
\newtcolorbox{exemplebox}[1][]{boxbase,colback=excol!5,colframe=excol,
  title={\textsc{Exemple}\ifx\relax#1\relax\relax\else\textmd{ : \itshape #1}\fi}}
\newtcolorbox{piege}[1][]{boxbase,colback=attncol!6,colframe=attncol,
  title={\textsc{Piège}\ifx\relax#1\relax\relax\else\textmd{ : \itshape #1}\fi}}
\newtcolorbox{remarque}[1][]{boxbase,colback=black!4,colframe=black!45,
  title={\textsc{Remarque}\ifx\relax#1\relax\relax\else\textmd{ : \itshape #1}\fi}}
% --- boîtes EXERCICE / CORRIGÉ (noms EXACTS, auto-comptées) ---
\newtcolorbox[auto counter]{exo}[1][]{boxbase,colback=primary!4,colframe=primary,
  title={\textsc{Exercice}~\thetcbcounter\ifx\relax#1\relax\relax\else\textmd{ --- \itshape #1}\fi}}
\newtcolorbox[auto counter]{corrige}[1][]{boxbase,colback=excol!4,colframe=excol,
  title={\textsc{Corrigé}~\thetcbcounter\ifx\relax#1\relax\relax\else\textmd{ --- \itshape #1}\fi}}
\newcommand{\vv}[1]{\vec{#1}}\newcommand{\ds}{\displaystyle}
\newcommand{\R}{\mathbb{R}}\newcommand{\N}{\mathbb{N}}\newcommand{\Z}{\mathbb{Z}}
% (un \usepackage{fancyhdr}+en-tête est possible ; il est ignoré côté web)
% =========================================================================
\begin{document}

\begin{center}
{\large\bfseries\color{excol} Thème 4 — Corrigé Moyen}\\[-2pt]
{\color{excol}\rule{5cm}{1pt}}
\end{center}

\begin{corrige}[fentes de Young : interfrange et frange]
On convertit en SI : $\lambda=\num{5e-7}\,\si{\metre}$, $a=\num{1e-3}\,\si{\metre}$.
\begin{enumerate}[label=\alph*)]
  \item $i=\dfrac{\lambda D}{a}=\dfrac{\num{5e-7}\times 3}{\num{1e-3}}
        =\num{1.5e-3}\,\si{\metre}=\SI{1.5}{\milli\metre}.$
  \item $x_2=2\,i=2\times 1{,}5=\SI{3}{\milli\metre}.$
\end{enumerate}
\end{corrige}

\begin{corrige}[nature de l'interférence en un point]
On calcule $\delta=|d_1-d_2|$ puis on compare à $\lambda=\SI{3}{\centi\metre}$.
\begin{enumerate}[label=\alph*)]
  \item $\delta=15-12=\SI{3}{\centi\metre}=\lambda$ (soit $\delta/\lambda=1$, entier) :
        interférence \textbf{constructive}.
  \item $\delta=13{,}5-12=\SI{1.5}{\centi\metre}=\lambda/2$ (soit $\delta/\lambda=0{,}5$) :
        interférence \textbf{destructive}.
\end{enumerate}
\end{corrige}

\begin{corrige}[remonter à la longueur d'onde]
\begin{enumerate}[label=\alph*)]
  \item $i=\dfrac{\lambda D}{a}\ \Rightarrow\ \lambda=\dfrac{i\,a}{D}.$
  \item $\lambda=\dfrac{\num{3e-3}\times\num{5e-4}}{2{,}5}
        =\dfrac{\num{1.5e-6}}{2{,}5}=\num{6e-7}\,\si{\metre}=\SI{600}{\nano\metre}.$
\end{enumerate}
\end{corrige}

\begin{corrige}[du déphasage à la différence de marche]
On utilise $\delta=\dfrac{\lambda}{2\pi}\,\Delta\varphi$.
\begin{enumerate}[label=\alph*)]
  \item $\Delta\varphi=\pi$ : \textbf{destructive}, $\delta=\lambda/2$.
  \item $\Delta\varphi=2\pi$ : \textbf{constructive}, $\delta=\lambda$.
  \item $\Delta\varphi=\pi/2$ : interférence \textbf{intermédiaire} (ni maximale, ni nulle),
        $\delta=\lambda/4$.
\end{enumerate}
\end{corrige}

\begin{corrige}[éloigner l'écran]
$i=\lambda D/a$ avec $\lambda=\num{6e-7}\,\si{\metre}$ et $a=\num{4e-4}\,\si{\metre}$.
\begin{enumerate}[label=\alph*)]
  \item $i=\dfrac{\num{6e-7}\times 1}{\num{4e-4}}=\num{1.5e-3}\,\si{\metre}=\SI{1.5}{\milli\metre}.$
  \item $i=\dfrac{\num{6e-7}\times 2}{\num{4e-4}}=\num{3e-3}\,\si{\metre}=\SI{3}{\milli\metre}.$
  \item $i$ est proportionnel à $D$ : doubler $D$ \textbf{double} l'interfrange.
\end{enumerate}
\end{corrige}

\end{document}
